\section{Main results}\label{sec:main}
Let us now come to the main objects in this work. We denote by $\mathbb{D}(\eta,r)$ the disc with centre $\eta \in \C$ and radius $r$. For a given sequence of points $p_N$, the positive number $r_N=r_N(p_N)$ is called the \textit{micro-scale} if it satisfies
\begin{gather} \label{micro-scale}
	\int_{ \mathbb{D}(p_N,r_N) } \frac{ \Delta Q(\zeta) }{2} \,{\rm d}A(\zeta)= \frac{1}{N}.
\end{gather}
We drop the subscript and write $p \equiv p_N$ if the sequence does not depend on $N$.
By \eqref{density}, the micro-scale $r_N$ corresponds to the mean eigenvalue spacing in radial distance of the ensemble~\eqref{Gibbs} at the point $p\in S$
(cf.~\cite{BE} for a situation where $p$ is outside the droplet).
We define the rescaled process $\boldsymbol{z}=\{ z_j \}_{ j=1 }^N$ as follows: for all $j$,
\begin{gather} \label{rescaling}
	z_j:=
	\begin{cases}
		r_N^{-1} \cdot (\zeta_j-p) &\text{if } p \in \operatorname{int}(S) ,
		\\
		{\rm e}^{-{\rm i}\theta} r_N^{-1} \cdot (\zeta_j-p) &\text{if } p \in \partial S ,
	\end{cases}
\end{gather}
distinguishing the interior and the boundary of the droplet $S$.
Here the angle $\theta\equiv \theta_N \in \R$ is chosen so that ${\rm e}^{{\rm i}\theta}$ is outer normal to $\partial S$ at $p$, see Figure~\ref{Fig. QGinibre}(b)
inset
for an illustration of the rescaled process.
The $k$-point correlation function of
the rescaled process $\boldsymbol{z}$ is defined by
\begin{gather*}
R_{N,k}(z_1,\dots,z_k)
:= \lim_{\eps \downarrow 0} \frac{\mathbb{P}(\exists \text{ at least one particle in } \mathbb{D}(z_j,\eps),\, j=1,\dots,k ) }{\eps^{2k}}.
\end{gather*}
In \eqref{bfRNk def} a more standard definition of the $k$-point correlation function before rescaling denoted by $\bfR_{N,k}(\zeta_1,\dots,\zeta_k)$ is given.
Throughout the paper we distinguish these and all other objects (kernels, Hamiltonian, joint distribution) before rescaling by bold symbols.
For the precise relation on the level of the kernels see~\eqref{kappa_N}. It results into the following relation between the $k$-point correlation functions
\begin{gather*}
R_{N,k}(z_1,\dots, z_k)=r_N^{2k} \bfR_{N,k}(\zeta_1,\dots,\zeta_k).
\end{gather*}
It is well known that before and after rescaling the
set $\boldsymbol{z}$ forms a Pfaffian point process (see \cite{MR1928853,Mehta}), i.e., $R_{N,k}$ is expressed in terms of a certain $2\times 2$ matrix-valued kernel $K_N$ as
\begin{gather} \label{RNk Pf}
	R_{N,k}(z_1,\dots, z_k)=\prod_{j=1}^{k}
(\bar{z}_j-z_j) \Pf \big[ K_{N}(z_j,z_l) \big]_{ j,l=1 }^k ,
\end{gather}
where the kernel $K_N$ is of the form
\begin{gather*}
	K_N(z,w):={\rm e}^{ -\frac{N}{2} (Q( \zeta )+Q( \eta )) } \begin{pmatrix}
		\kappa_N(z,w) & \kappa_N(z,\bar{w}) \\
		\kappa_N(\bar{z},w) & \kappa_N(\bar{z},\bar{w})
	\end{pmatrix}.
\end{gather*}
Here $\Pf$ denotes a Pfaffian of the $2k \times 2k$ skew-symmetric matrix and
\begin{gather} \label{zeta eta}
\zeta=\begin{cases}
	p+r_N z &\text{if }p \in \textup{int}(S),
	\\
	p+{\rm e}^{{\rm i}\theta} r_N z &\text{if }p \in \partial S,
\end{cases} \qquad
 \eta=\begin{cases}
	 p+r_N w &\text{if }p \in \textup{int}(S),
	\\
	 p+{\rm e}^{{\rm i}\theta} r_N w &\text{if }p \in \partial S,
\end{cases}
\end{gather}
where $\theta$ is given as in \eqref{rescaling}.
The arguments \eqref{zeta eta} are given according to the rescaling \eqref{rescaling}.
For the spectral density at $k=1$, let us denote $R_N \equiv R_{N,1}$.

